

\chapter{Introduction}

Images captured in high-altitude or extreme-weather regions such as fog, smog, dust storms, or snow often suffer from low contrast and color distortion.  
These degradations create a distribution shift between clean datasets and real-world imagery, causing transformer-based vision systems like ViT~\cite{dosovitskiy2020vit} and multimodal encoders such as CLIP~\cite{radford2021clip} to perform unreliably.  
Most modern vision backbones, derived from the transformer paradigm~\cite{vaswani2017attention}, assume clear, well-lit scenes and therefore fail when deployed on aerial or drone cameras in low-visibility conditions.
This project addresses that gap by generating synthetic data that emulate such adverse environments.  
Rather than relying on scarce real samples, diffusion-based generative models recreate realistic images exhibiting fog, cloud cover, pollution-induced smog, or heavy snow.  
By exposing vision networks to these synthetically shifted distributions, we aim to improve robustness and interpretability under natural domain variations, aligning with the broader generative learning frameworks that evolved from VAEs~\cite{kingma2013autoencoding} and GANs~\cite{goodfellow2014gan}.
Text and image-conditioned diffusion frameworks control the degree and type of degradation, enabling the creation of high-fidelity yet distribution-shifted datasets that support vision–language systems used in drone-based agricultural surveys at high-altitude.

\section{Problem Statement}


Current vision systems rarely encounter extreme environmental conditions and thus fail to generalize beyond clean training domains.
This work focuses on \emph{dataset synthesis} using diffusion models to generate realistic, distribution-shifted images under fog, mist, and snow.
The goal is to design a controllable generation pipeline that yields diverse data and enables adaptation of pre-trained models to these scenarios, enhancing perception reliability in adverse weather.
Collecting real high-altitude imagery is difficult and unsafe, while existing datasets remain scarce and unrepresentative.
Generative modeling offers a practical way to replicate these conditions synthetically, producing realistic yet diverse images that preserve structure while incorporating degradations typical of extreme weather crucial for reliable performance in drone-based applications.
